\section{Introduction}

We trained the same code language model on four different practice problem sets and found that
practicing HumanEval-style algorithm completions made the model better at MBPP utility scripts
than practicing MBPP scripts directly --- by up to 1.5 percentage points, consistently across
all three random seeds.

Code LLMs are routinely fine-tuned on mixes of publicly available datasets --- HumanEval
training problems, MBPP utility scripts, LeetCode algorithmic challenges --- but practitioners
choose these sources based on intuition rather than measurement. The prevailing assumption is
that training on X specializes a model for X: MBPP training should confer MBPP benchmark
advantage; HumanEval training should sharpen HumanEval performance. If this intuition were
correct, source selection would be straightforward. It is not correct.

The deeper problem is that no controlled study has isolated the causal effect of SFT source
identity on code benchmark pass@1 while holding model architecture, token budget, training
hyperparameters, and supervision format constant. Existing code SFT papers vary data quality,
quantity, and format simultaneously, making causal attribution impossible. Domain mixture
optimization literature \cite{Xie2023DoReMi,Zhang2026DomainPilot} addresses pretraining but has not
been applied as a source-identity ablation for code-specific SFT with execution-based
evaluation. The gap --- a controlled 4-source $\times$ 2-benchmark transfer matrix with token-budget
equalization --- does not exist in the literature.

We close this gap and discover that source selection is not a neutral engineering choice: it
is the dominant factor governing post-SFT benchmark performance. Under identical token budgets,
HumanEval-only training achieves 35.0\% HumanEval+ pass@1 while LeetCode-only training achieves
$\sim$3.1\% --- a $\sim$32 percentage point gap for the same model on the same compute budget. This effect
substantially exceeds most reported architectural improvements.

The key insight enabling our analysis: \textbf{code-embedding cosine similarity between the SFT
training source and the test benchmark perfectly predicts the pass@1 rank order across all four
conditions.} When we rank the four SFT sources (HumanEval-only, MBPP-only, Equal-mix,
LeetCode-only) by their CodeBERT embedding similarity to HumanEval+, the order is
HE $>$ MB $>$ EQ $>$ LC. This matches the observed HumanEval+ pass@1 rank exactly (Spearman
$\rho$=1.0, permutation test p=0.042, n=10,000 shuffles). Embedding-space distributional alignment
is a measurable, pre-training predictor of SFT behavioral outcomes.

This insight also explains the counterintuitive finding: HumanEval-style algorithmic function
completion problems are embedding-closer to both HumanEval+ and MBPP+ than MBPP utility
scripts are. The training distribution that best aligns with the test distribution ---
regardless of surface-level similarity --- produces the best performance.

Building on this insight, we make the following contributions:

\begin{enumerate}
\item \textbf{First controlled source-identity ablation for code SFT}: We train DeepSeek-Coder-1.3B/7B-Base
  on four isolated conditions (HumanEval-only, MBPP-only, LeetCode-only, Equal-mix) with
  token-budget equalization, deduplication, and format normalization --- enabling causal
  attribution of the 31.9pp source identity effect (ANOVA F=11.37, p=0.020).

\item \textbf{Distributional alignment predicts code SFT rank order}: CodeBERT embedding cosine
  similarity between training source and test benchmark perfectly predicts HumanEval+ pass@1
  rank across all four conditions ($\rho$=1.0, p=0.042, dual-encoder concordant), providing the
  first permutation-tested evidence that embedding-space alignment governs code SFT
  specialization.

\item \textbf{HumanEval SFT confers cross-benchmark coding advantage}: Contrary to symmetric
  specialization, HumanEval-only training achieves the highest pass@1 on both HumanEval+
  (35.0\%) and MBPP+ ($\sim$52\%), suggesting that algorithmic function-completion training develops
  more transferable Python programming competence than utility-script training.

\item \textbf{Methodological insight on scale comparison}: At 7B scale, within-condition seed variance
  collapses to near-deterministic ($\sigma^2$=0.00043), inflating $\eta^2$ and motivating absolute
  condition spread as the primary scale-comparison metric.
\end{enumerate}

We organize the paper as follows: Section~\ref{sec:related} surveys related work.
Section~\ref{sec:method} describes our methodology.
Section~\ref{sec:setup} presents the experimental setup.
Section~\ref{sec:results} reports results.
Section~\ref{sec:discussion} discusses findings and limitations.
Section~\ref{sec:conclusion} concludes.
